
\section{候选方向比较}
在完成选题的重新定义之后，接下来的问题是在可选的研究方向中，哪一条最适合在当前时间、已有系统基础和答辩要求下收敛成一个成立的毕设课题。从目前的讨论来看，真正有代表性的方向大致可以分成四类：继续沿蛋白质设计系统增强路线推进、转向自动benchmark、转向工作流编译/验证/修复，以及转向运行时 controller。它们分别对应不同的研究对象、不同的创新重心，也对应不同的风险。

\subsection{延续“蛋白质设计系统增强”路线}

这是与当前系统最接近的一条路线。其基本思路是继续围绕生命科学任务展开，将系统定义为一个面向蛋白质设计或更广义生命科学任务的多Agent助手，并再次基础上不断补充能力，例如更完整的工作链、更复杂的任务拆解、更细致的门控、更丰富的HITL交互，以及更完善的恢复逻辑。

这一条路线的核心是：把系统当作应用产品或综合平台来完善。其关注重点在于系统覆盖了多少能力、能处理多少真实请求、交互链路是否顺畅、执行流程是否稳定。它的优势十分明显：与现有系统最贴近，实现成本最低，也最容易展示出成果。

但它的问题同样明显。首先，这条线路最容易被看成是另一个生命科学助手的自然扩展。因为只要系统的对外叙事仍然是“面向生命科学任务（蛋白质设计）的多Agent帮助系统”，那么内部无论加入多少状态机、恢复机制和人工确认逻辑，本质上都没有脱离“助手系统”这一概念。其次，这一条路径的创新点很难集中，往往会分散在多个工程增强点上，例如增加某一类工具、改进某一段流程或补充某种门控逻辑。这些内容单独来看都合理，但是很难压缩成一个明确的算法主线。最后，在当前时间压力下，继续走“大系统增强”的路线很容易导致内容越做越多，主线却越来越模糊。

因此，虽然这条路线工程上最平滑，但是从毕设选题的角度来看，它恰恰是风险最高的一条，因为它最容易落入无法解释核心研究问题是什么的局面。

\subsection{自动Benchmark路线}

自动benchmark路线的思路是把研究对象从“如何更好地完成任务”转向“如何系统地评测复杂科学工作流任务”。在这一方向上，系统不再首先关注“输出什么结果”，而是关注：如何构造具有不同复杂度和难度层级的任务，如何让这些任务可复现、可比较、可自动评分，以及如何让benchmark真正反映多Agent系统在高代价科学工作任务中的能力差异。

从近两年的研究趋势来看，这条路线是成立的，而且学术价值很明确。PaperBench\parencite{paperbench}强调的是研究复现实验的细粒度自动评分与rubric拆解，BrowseComp\parencite{browsecomp}强调的是高难度信息寻求任务中的持久搜索，CoLLAB\parencite{collab}则突出协作复杂度和多Agent benchmark的可扩展构造。

这一线路的本质是：把“难度”本身变成研究对象。它不像助手系统那样关注“给出什么答案”，而是关注“什么样的任务更难、为什么难、如何被自动构造出来、又如何被可靠评分。”从独立性上来说，这条线路非常安全，因为benchmark不是生命科学助手的子模块，反而可以把不同助手都作为benchmark的被评测对象。从答辩角度看，也更容易回答“贡献到底是什么”，因为benchmark一旦设计得好，本身就是独立研究资产。

不过，这条线路本身也有自己的问题。首先，它的重点不是控制算法，而是在任务生成、难度分层和自动评分机制上。如果自动生成和自动评分不够强，benchmark很容易退化成“整理了一批测试集”。其次，这条线路虽然能很好地体现“难度”，但是也相对弱化了与当前系统直接结合的优势。也就是说，它更独立，但是与现有系统的耦合不够。最后，从剩余时间来看，如果要同时做好任务构造、难度校准和轨迹级评分，工作量其实不小，切容易在“设计评测体系”这一块耗费过多时间。

因此，自动benchmark是一条边界很清晰、不容易成为子集的路线，但是它更像是“独立评测平台”选题，和当前系统的天然结合度略低。

\subsection{工作流编译、验证与修复路线}

另一条可能的方向，是把研究结果定义为“从科研目标到可执行工作流的编译、验证与修复”。在这条路线下，重点不是任务本身，也不是最终结果，而是：如何从自然语言目标生成结构化工作流；如何检查工具链的I/O兼容性、参数合法性和资源约束；以及如何在执行失败后进行最小修复和后缀重规划。

从方法论上来看，这条路线非常有计算机专业特色，因为它天然带有\texttt{compiler}、\texttt{verification}、\texttt{schema caching}、\texttt{plan repair}等色彩。现有多\texttt{Agent}系统设计中，越来越多工作开始强调结构化任务表示、明确的状态机组织以及对工具链和交互协议的形式约束。例如MetaAgent\parencite{metaagent}这类工作已经开始从显式状态机组织的角度构造多Agent系统，而工作流系统本身也越来越强调plan schema、tool contract和可验证执行过程。

这条线路的核心是：把科学任务的执行链看作一种可编译、可验证、可修复的程序结构。如果从理论表达上看，它很有吸引力。

但是它在当前语境下有两个显示问题。第一，与现有系统已有能力重合较多。当前系统已经具有显式FSM、计划候选、patch/replan、候选验证和恢复机制，因此如果再把“工作流编译与修复”作为核心方向，很容易被看成对现有\texttt{planner}的进一步完善，而不是一个研究主线。第二，这条路线虽然很具有形式化，但是在叙事上较弱。换句话说，它更像“让流程更正确”，而不是“让系统在高代价下做出更好的运行时决策”。从应用意义和系统价值上，可能不够直观。

因此，工作流编译/验证/修复路线是一个理论性较强的候选方向，但是与当前系统的差异化空间并没有想象中大，而且更容易被视为现有\texttt{planner}和恢复逻辑的深化版本。


\subsection{运行时Controller路线}

\texttt{controller}路线的出发点是把研究对象上升为：在高代价、长链路、带恢复与人工介入的科学工作流中，系统究竟如何决定下一步如何推进。在这一方向上，系统不再被理解成一个“完成生命科学任务的助手”，而被理解成一个“治理工作执行的运行时控制器”。

从近两年的研究来看，这条方向虽然还没有完全同构的现成答案，但是已经有多条重要脉络构成了它的理论背景。在不稳定性建模方面，不确定性思维（Uncertainty of Thoughts，UoT）说明不确定性可以进入\texttt{Agent}的规划闭环\parencite{uot}；在推理资源方面，BEST-Route 方法提出依据任务难度与质量阈值，对推理计算资源进行动态调度，实现性能与成本之间的权衡\parencite{best-route}；在推理过程优化方面，QLASS 方法通过对中间决策状态进行价值估计，引导 Agent 在推理空间中的搜索路径，提高推理效率与结果质量\parencite{qlass}；在行动与工具使用层面，EcoAct 从经济性约束出发，对 Agent 的动作选择与工具调用策略进行建模，强调资源受限环境下的最优决策问题\parencite{ecoact}；在多 Agent 系统层面，Why Do Multi-Agent LLM Systems Fail? 与 Which Agent Causes Task Failures and When? 的研究表明，系统失败并不完全来源于单一模型能力的不足，而与协作机制、验证流程、责任归因以及系统治理结构密切相关\parencite{whyfail,whichandwhen}。

这一路线本质是：把“难度、风险、工作量和人工介入价值”统一纳入到一个运行时的决策问题。换句话说，\texttt{controller}关注的不是任务内容本身，而是系统治理动作本身。它所决定的不是“下一句怎么回答”，而是诸如：

\begin{itemize}
	\item 现在是否值得继续执行；
	\item 是否应该扩大候选探索
	\item 是否应切换到保守路径；
	\item 当前失败更适合 \texttt{patch}还是\texttt{replan}；
	\item 是否应请求\texttt{HITL}
	\item 是否应在进一步浪费成本前提前终止
\end{itemize}

这使得\texttt{controller}方向同时具有三方面优势。第一，它与当前系统切合度最高。因为现有系统已经有状态机、HITL、patch/replan、阶段化执行链和实验框架，\texttt{controller}可以自然地嵌入这些现有能力之上，而不需要推翻架构。第二，它最容易把“难度”从口头概念变成真正的算法变量。难度不再指示任务描述中的印象，而是需要在执行中动态估计、并据此改变策略的核心状态。第三，它在应用意义上最容易成立。蛋白质设计这类高代价场景天然就会遇到“什么时候继续算、什么时候该求助人类、什么时候必须止损”的问题，因此\texttt{controller}的价值不仅是理论上的，也非常符合真实科研流程的痛点。

当然，\texttt{controller}方向也不是没有风险。它最大的风险在于，如果最后只做成“几个特征+一堆人工阈值+超阈值就进入HITL”的规则系统，那么它虽然还能运行，但创新性会明显不足。因此这条路线真正成立的前提是：\texttt{controlle}必须被明确建模为一个部分可观测、带风险与预算约束的序贯决策问题，而不是一个更复杂的门控模块。

\subsection{不同方向的比较与理由}

综合来看，这几个方向的差异可以概括为：生命科学助手最贴近现有系统，但独立性最弱，也最容易被已有父集覆盖；自动benchmark路线独立性最强，难度问题也最自然，但与现有系统直接结合较弱，且任务生成与自动评分部分在剩余时间内工作量也不小；工作流编译/验证/修复路线理论上较“计算机化”，但与现有\texttt{planner}的功能边界较近，容易被视为现有平台能力的深化；而\texttt{controller}路线则位于一个相对平衡的位置，它既能明显提升选题层次，使课题从“助手系统”转向“工作流治理”，又能最大限度复用现有系统作为实验平台，同时还能把“难度、风险、工作量、恢复和HITL”这些此前分散的概念统一到一个核心决策问题中。

因此，如果在当前时间节点下选择一条既能体现算法性，又能形成明确主线、同时还能落到现有系统的路径，那么controller是最值得优先考虑的方向。而在controller内部，我准备进一步倾向于信念状态（Belief State)controller，因为它最适合把“难度”定义为一个运行时的隐变量，而不是一个静态标签或单纯的打分结果。


